%! Introducing Statistics: What Can We Learn from Data? — Unit 1, Lesson 1.0 — Experience & Formalize
%! (Activity · QuickNotes · Application · Check Your Understanding) (Answer Key)
% Mirrors the blank component byte-for-byte apart from the package swap, the pageheader, and
% the filled answers. The \answerspace macro is IDENTICAL to the blank's — that identity is
% what keeps this key page-for-page with the blank.
% NO teachernote here — teacher prose lives in the lesson plan.
\documentclass[12pt]{article}
\usepackage{apstats-article}
\usepackage{apstats-key}   % pulls in -boxes; do NOT also load -boxes
\newcommand{\answerspace}[2]{\par\nopagebreak\noindent\begin{minipage}[t][#1][t]{\linewidth}%
  \color{keyred}\bfseries #2\end{minipage}\par}

\begin{document}

\pageheader{Unit 1, Lesson 1.0}{Experience \& Formalize: The Adoption Board --- Answer Key}
% NAMESTRIP: no \namedateperiod here — mirrors the blank.

\vspace{0.06in}

\begin{headlinebox}{sky}
\textbf{The Adoption Board.} A county animal shelter posts one board in its lobby listing every
dog currently up for adoption. Visitors read it, and so does the shelter director --- but they
are not looking for the same things. With your group, work through the two problems below using
only what you already know, and be ready to explain \emph{how you know} each answer
\textbf{from the board itself}.
\end{headlinebox}

\vspace{0.04in}

\begin{scenariobox}[1. Reading the Board]{navy}
\small
\renewcommand{\arraystretch}{1.25}
\begin{tabularx}{\linewidth}{@{} l Y Y Y Y @{}}
\rowcolor{sky}
\textbf{Name} & \textbf{Breed} & \textbf{Age (yr)} & \textbf{Weight (lb)} & \textbf{Days on board} \\
\hline
Biscuit & Beagle    & 3 & 24 & 12 \\
Momo    & Terrier   & 1 & 15 & 44 \\
Ranger  & Shepherd  & 5 & 61 &  6 \\
Pepper  & Retriever & 2 & 55 & 21 \\
Sage    & Beagle    & 7 & 27 & 59 \\
Tuck    & Mixed     & 4 & 38 &  9 \\
Willow  & Terrier   & 2 & 17 & 33 \\
Zeke    & Shepherd  & 6 & 70 & 15 \\
\hline
\end{tabularx}

\normalsize
\begin{enumerate}[label=\textbf{\alph*.}, leftmargin=1.8em, itemsep=8pt]
  \item One \textbf{row} of this board describes \ans{one dog}.

        The board has \ans{8} rows and \ans{4} columns of recorded facts
        (not counting the Name column).
  \item ``Beagle'' appears on the board. Is it a \textbf{column heading}, or something
        \textbf{written underneath} one? Circle your answer, then say which column heading it
        belongs to. \hfill \textbf{heading} \qquad \ans{\textbf{written underneath}}
        \answerspace{1.4cm}{It is written underneath the heading \emph{Breed}. ``Breed'' is
        what was recorded; ``Beagle'' is one of the answers recorded under it.}
  \item \textbf{Underline} the part of the board you would read to answer each question.
        One needs a single spot; the other needs a whole column.

        \emph{How much does Zeke weigh?} \qquad \emph{How much do these dogs weigh?}
        \answerspace{1.6cm}{Zeke: one cell --- 70 lb. All the dogs: the whole Weight column
        (24, 15, 61, 55, 27, 38, 17, 70), because there is no single answer.}
  \item Sage has been on the board for 59 days and Ranger for 6. If \emph{every} dog had been
        on the board the same number of days, would the ``Days'' column still be worth posting?
        Explain.
        \answerspace{2.0cm}{No. If every entry were the same, one dog would tell you everything
        about waiting time and the column would carry no new information. The column is worth
        posting precisely \emph{because} the waits differ from dog to dog.}
\end{enumerate}
\end{scenariobox}

\boxguard[16]
\begin{scenariobox}[2. Four Questions the Director Might Ask]{navy}
The shelter director looks at the same board and asks these four questions.

\begin{enumerate}[label=\textbf{\alph*.}, leftmargin=1.8em, itemsep=8pt]
  \item For each one, decide whether you could answer it by reading \textbf{one spot} on the
        board, or whether you would have to read \textbf{many rows}. Write \emph{one spot} or
        \emph{many rows} in the blank.

        \begin{enumerate}[label=\textbf{(\Alph*)}, leftmargin=2.2em, itemsep=5pt]
          \item Is Momo a terrier? \hfill \ans{one spot}
          \item How many dogs are on the board? \hfill \ans{many rows}
          \item How long do dogs on this board tend to wait? \hfill \ans{many rows}
          \item Do heavier dogs tend to wait longer than lighter ones? \hfill \ans{many rows}
        \end{enumerate}
  \item Two of those needed many rows. But only \emph{one} of the two has an answer that
        \textbf{changes from one dog to the next}. Which one is which, and how can you tell?
        \answerspace{2.6cm}{(B) needs many rows but has a single fixed answer --- 8 --- that
        does not belong to any one dog, so nothing varies. (C) has a different answer for every
        dog (12, 44, 6, 21, \dots), so its answer varies across individuals. Counting the rows
        is not the same as measuring something that differs between them.}
  \item A flyer taped under the board reads: \textbf{``Our dogs go home in 25 days.''}
        Name \textbf{two} things the flyer does not tell you that you would want to know
        before believing it.
        \answerspace{2.0cm}{Any two of: which dogs it counted and when; how many dogs;
        what ``25 days'' summarizes (a typical value? the longest?); whether it includes dogs
        still waiting; how much the waits spread out.}
  \item A visitor reads the flyer and says, ``Good --- so I'll have my dog home in 25 days.''
        Using the board, give one specific reason that is the wrong way to read the number.
        \answerspace{2.0cm}{The board's waits run from 6 days (Ranger) to 59 days (Sage), so
        no single dog is guaranteed 25. A number that summarizes a group does not predict any
        one individual in it.}
\end{enumerate}
\end{scenariobox}

\boxguard[14]
\begin{tcolorbox}[enhanced, colback=sky, colframe=navy, boxrule=0.8pt, arc=2mm,
    left=3mm, right=3mm, top=1.5mm, bottom=1.5mm, breakable,
    fonttitle=\bfseries\small\color{white}, title={QuickNotes: Data, Variables, and Statistical Questions},
    attach boxed title to top left={yshift=-2mm, xshift=4mm},
    boxed title style={colback=navy, colframe=navy, arc=1.5mm}]
\small
\begin{minipage}[t]{0.36\linewidth}
\vspace{2pt}
\renewcommand{\arraystretch}{1.15}
\begin{tabular}{@{}l l@{}}
\multicolumn{2}{@{}l}{\footnotesize\textbf{\color{navy}Breed} \quad \textbf{\color{navy}Weight}}\\
\hline
Beagle    & 24 \\
Terrier   & 15 \\
Shepherd  & 61 \\
\hline
\end{tabular}\\[3pt]
{\footnotesize Headings are \textbf{variables}.\\
What is written under them\\
are \textbf{values}. Each row is\\
one \textbf{individual}.}
\end{minipage}\hfill
\begin{minipage}[t]{0.60\linewidth}
\vspace{-4pt}
\begin{itemize}[leftmargin=1.1em, itemsep=5pt]
  \item \textbf{Data} are recorded values \emph{plus} the \ans{context} that gives them meaning.
        A number alone answers nothing.
  \item Before a number is data, answer the \textbf{W questions}: who, \ans{what},
        in what \ans{units}, when and where.
  \item An \textbf{individual} is what one \ans{row} describes. A \textbf{variable} is a
        \ans{column} heading --- never one of the entries under it.
  \item \textbf{Variability:} values \ans{differ} from one individual to the next.
        \emph{This is why statistics exists.}
  \item A \textbf{statistical question} passes \textbf{both} checks:\\[2pt]
        \textbf{1.} Do the answers \ans{vary} across individuals?\\[2pt]
        \textbf{2.} Must I \ans{collect} data to answer it?
\end{itemize}
\end{minipage}
\end{tcolorbox}

\boxguard[14]
\begin{notesbox}{Application: The Gym Sign-Up}
We will do this one together. A gym records the \textbf{resting heart rate}, in beats per
minute, of each of its 120 members on the morning they join. It also records each member's
\textbf{home ZIP code}.

\begin{enumerate}[label=\textbf{\alph*.}, leftmargin=1.8em, itemsep=7pt]
  \item The individuals are \ans{the 120 gym members}, and two variables recorded are
        \ans{resting heart rate} and \ans{home ZIP code}.
  \item Write one statistical question this gym could answer, and say exactly what varies in it.
        \answerspace{1.8cm}{``How do the resting heart rates of the gym's members vary?''
        What varies is the resting heart rate, in bpm, from one member to the next.}
  \item \textbf{What if we change the variable?} ``What is the most common ZIP code among the
        members?'' --- is that a statistical question? Run both checks before you answer.
        \answerspace{2.2cm}{Yes. Check 1: members live in different ZIP codes, so the answer
        differs from member to member. Check 2: you must collect the data to know. A variable
        does not have to be a number --- ZIP code is a label, and it still varies.}
\end{enumerate}
\end{notesbox}

\boxguard[14]
\begin{notesbox}{Check Your Understanding}
Work with your partner, then finish on your own. \emph{These are for practice --- they are not
scored.}
\begin{enumerate}[leftmargin=1.9em, itemsep=9pt]
  \item A biologist weighs each of the 30 turtles in a pond and records the weight in kilograms.

        The individuals are \ans{the 30 turtles}; the variable is \ans{weight},
        measured in \ans{kilograms}.
  \item Write one statistical question about those turtles, then name what varies in it.
        \answerspace{1.8cm}{``How do the weights of the turtles in this pond vary?''
        What varies is weight, in kg, from one turtle to the next.}
  \item \textbf{Same pond, two questions.} Explain why exactly one of these is statistical.

        \emph{(i)} How much does the heaviest turtle weigh? \quad
        \emph{(ii)} How do the turtles' weights vary?
        \answerspace{2.0cm}{(ii) is statistical. (i) has one fixed answer describing a single
        turtle, so nothing varies; (ii) has a different answer for each turtle and can only be
        answered by collecting all 30 weights.}
  \item A report states only: \textbf{``The average is 12.''} Name \textbf{two} pieces of
        context you would need before this number means anything.
        \answerspace{1.8cm}{Any two of: who or what was measured, what characteristic was
        measured, in what units, and when or where the data were collected.}
  \item \textbf{Formative check.} Circle the statistical question, then justify it in one
        sentence by pointing to what varies.
        \begin{enumerate}[label=\textbf{(\Alph*)}, leftmargin=2.2em, itemsep=4pt]
          \item How many students are enrolled at this school?
          \item How many hours did Jordan work last week?
          \item How do the hours worked each week vary among students who have jobs?
                \ans{$\leftarrow$ correct}
          \item What time does the library close on Friday?
        \end{enumerate}
        \answerspace{1.6cm}{(C): hours worked differ from one working student to the next, so
        the question has many answers and needs data. (A) and (D) each have one fixed answer;
        (B) describes a single individual.}
\end{enumerate}
\end{notesbox}

\end{document}
